% Copyright (c) 2026 Eugene Joseph M. Ragasa.
% SPDX-License-Identifier: Apache-2.0
% Author: Eugene Joseph M. Ragasa
% Formal proof completed: 2026-10-27
\documentclass[11pt]{article}
\usepackage[T1]{fontenc}\usepackage{amsmath,amssymb}\usepackage[margin=1in]{geometry}\usepackage[hidelinks]{hyperref}
\title{Proof Dossier: Trivial Scalar Operator Conjugation}\author{Eugene Joseph M. Ragasa}\date{Formal proof completed: 2026-10-27}
\begin{document}\maketitle
\section{Identity and status}
ID \texttt{BG-FOURIER-THM-049}; status \texttt{proved}. Lean declaration:
\par\noindent{\tiny\path{ProjectKoios.BandGauge.Fourier.}}\\
{\tiny\path{scalarPhaseConjugation_eq}}.
\section{Statement and proof}
Conjugating a complex scalar operator by any momentum-dependent unit-modulus phase leaves
it unchanged pointwise. Commutativity moves the inverse phase next to the phase, whose
product is one.
\section{Boundary}
Rank-one scalar conjugation is operator-trivial and is not evidence of gauge-dependent
locality. Kernel axioms are \texttt{propext}, \texttt{Classical.choice}, and
\texttt{Quot.sound}; no admissions. Checked with Lean \texttt{4.35.0} and Mathlib
revision \texttt{56a005855ea17e0b0958f73ec801ae140b0a1679}. AI-assisted and
kernel-checked; no independent expert review recorded.
\end{document}
